\documentclass[11pt, a4paper]{article}
\usepackage[utf8]{inputenc}
\usepackage[T1]{fontenc}
\usepackage[brazilian]{babel}
\usepackage{lmodern}
\usepackage{geometry}
\usepackage{array}
\usepackage{longtable}
\usepackage{booktabs}
\usepackage{enumitem}

\geometry{a4paper, left=2.0cm, right=2.0cm, top=1.8cm, bottom=1.8cm}
\setlength{\parindent}{0pt}
\setlength{\parskip}{0.18cm}
\renewcommand{\arraystretch}{1.28}

\begin{document}

\begin{center}
    {\Large\bfseries ESPECIFICAÇÃO TEXTUAL DO MODELO DE CASOS DE USO}\\[0.15cm]
    {\large Sistema IQArMobi --- Monitoramento Inteligente da Qualidade do Ar em \emph{Smart Campus}}\\[0.15cm]
    {\small\textbf{Autor:} Noedson Romario Vieira da Silva \quad|\quad \textbf{PPGCA -- NDAE/UFPA}}
    \vspace{0.15cm}
    \hrule
\end{center}

\vspace{0.15cm}

Este documento apresenta a especificação textual detalhada dos casos de uso do \textbf{Sistema IQArMobi}, elaborada em estrita conformidade com o modelo estruturado apresentado na vídeo-aula da disciplina pelo Prof.\ Dr.\ Rodrigo Quites Reis (contemplando os campos: \emph{Identificador}, \emph{Nome}, \emph{Breve Descrição}, \emph{Atores e Participação}, \emph{Prioridade}, \emph{Pré-condições}, \emph{Pós-condições}, \emph{Fluxo de Eventos Principal}, \emph{Subfluxos}, \emph{Fluxos Secundários/Exceções} e \emph{Requisitos Não Funcionais Associados}).

% ============================================================================
% CASO DE USO 01 (Integrando UC01, UC02 e UC03 do Diagrama)
% ============================================================================
\section*{1. Especificação do Caso de Uso UC01}

\begin{longtable}{|>{\bfseries}p{3.8cm}|p{12.4cm}|}
\hline
\multicolumn{2}{|c|}{\textbf{UC01 --- Receber Telemetria Ambiental via WebHook}} \\ \hline

Identificador & \texttt{UC01} \\ \hline

Nome & Receber Telemetria Ambiental via WebHook \\ \hline

Breve Descrição & Este caso de uso permite que o sistema receba automaticamente os pacotes de telemetria dos sete compostos gasosos enviados pelo servidor de rede LoRaWAN, valide e armazene as leituras ambientais, acione obrigatoriamente a inferência de Inteligência Artificial (\texttt{UC02}) e, caso seja detectada anomalia física ou risco ambiental, dispare um alerta preventivo (\texttt{UC03}). \\ \hline

Atores & 
\begin{minipage}[t]{\linewidth}
\vspace{0.05cm}
\begin{tabular}{@{}p{4.2cm}p{7.8cm}@{}}
\textbf{Ator} & \textbf{Como participa do Caso de Uso} \\ \hline
\emph{Servidor Chirpstack} (Ator Primário / Sistema Externo) & Inicia o caso de uso enviando uma notificação HTTP (WebHook) contendo a identificação da estação sensora (\texttt{DevEUI}) e as concentrações decodificadas dos gases em ppm. \\
\end{tabular}
\vspace{0.1cm}
\end{minipage} \\ \hline

Casos de Uso Relacionados & 
\begin{itemize}[nosep, leftmargin=0.5cm]
    \item \textbf{\texttt{<<include>>} UC02 (Executar Inferência Multi-Task de Qualidade do Ar):} executado obrigatoriamente a cada nova leitura válida recebida.
    \item \textbf{\texttt{<<extend>>} UC03 (Emitir Alerta de Anomalia ou Violação):} executado opcionalmente quando o ponto de extensão \emph{``Anomalia ou Violação CONAMA/EPA Detectada''} é atingido após a inferência.
\end{itemize} \\ \hline

Prioridade & \textbf{Essencial (Alta)} \\ \hline

Pré-condições & 
\begin{enumerate}[nosep, leftmargin=0.5cm]
    \item A estação sensora emissora deve estar previamente cadastrada e ativa no sistema.
    \item O modelo neural \emph{Multi-Task} e os normalizadores devem estar carregados em memória pelo Agente Inteligente.
\end{enumerate} \\ \hline

Pós-condições & 
\begin{enumerate}[nosep, leftmargin=0.5cm]
    \item A nova leitura ambiental dos sete gases é registrada no histórico do sistema.
    \item A previsão futura (15 minutos à frente), a classificação do Índice de Qualidade do Ar (IQAr) e o diagnóstico de integridade do sensor são salvos e disponibilizados nos painéis.
    \item Se houver risco ambiental ou falha física do sensor, um registro de alerta ativo é gerado.
\end{enumerate} \\ \hline

Fluxo de Eventos Principal (Caminho Feliz) & 
\begin{enumerate}[nosep, leftmargin=0.5cm]
    \item O \textbf{Servidor Chirpstack} envia ao sistema um evento de telemetria (\texttt{up}) contendo o identificador da estação (\texttt{DevEUI}) e as medições dos sete gases (CO, CO\textsubscript{2}, NH\textsubscript{4}, álcool, acetona, tolueno e O\textsubscript{3}).
    \item O \textbf{Sistema} verifica a autenticidade do evento, identifica a estação sensora correspondente e registra a leitura ambiental no banco de dados.
    \item O \textbf{Sistema} agrega as leituras em janela temporal de 30 minutos (com médias de 1 minuto) e calcula a média móvel de 8 horas de monóxido de carbono (CO).
    \item O \textbf{Sistema} inclui e executa o caso de uso \textbf{UC02 (Executar Inferência Multi-Task de Qualidade do Ar)}, obtendo: (a) as concentrações previstas para os próximos 15 minutos, (b) a classificação normativa do IQAr (\emph{Boa}, \emph{Moderada} ou \emph{Ruim}) e (c) o diagnóstico de saúde física do sensor.
    \item O \textbf{Sistema} armazena os resultados da inferência vinculados à estação sensora.
    \item O \textbf{Sistema} retorna a confirmação de recebimento e processamento com sucesso (\texttt{HTTP 200 OK}) ao \textbf{Servidor Chirpstack}, encerrando o caso de uso.
\end{enumerate} \\ \hline

Fluxos Secundários (Alternativos e de Exceção) & 
\begin{itemize}[nosep, leftmargin=0.5cm]
    \item \textbf{[Fluxo Alternativo / Extensão no Passo 4 --- Violação Ambiental ou Anomalia no Sensor (\texttt{UC03})]:} Caso o resultado da inferência indique qualidade do ar futura \emph{``Ruim''} ($\ge 9{,}5\textrm{ ppm}$) ou falha física no sensor (erro de reconstrução elevado ou sinal congelado com variância nula), o \textbf{Sistema} aciona a extensão \textbf{UC03 (Emitir Alerta de Anomalia ou Violação)}, verifica se não há alerta idêntico emitido nos últimos 15 minutos (\emph{cooldown}) e registra um novo alerta crítico antes de prosseguir para o Passo 5.
    \item \textbf{[Exceção no Passo 2 --- Evento de Controle sem Carga de Telemetria]:} Caso o \textbf{Servidor Chirpstack} envie um pacote apenas de controle de rede (como confirmação de conexão \emph{join} ou status de bateria) sem leituras de gases, o \textbf{Sistema} ignora a etapa de inferência, registra apenas o status de comunicação e retorna confirmação ao servidor.
    \item \textbf{[Exceção no Passo 3 --- Histórico Insuficiente na Inicialização (\emph{Cold Start})]:} Caso a estação sensora tenha acabado de ser ligada e ainda possua menos de 30 minutos de histórico acumulado, o \textbf{Sistema} armazena a leitura atual, preenche temporariamente a janela histórica e sinaliza que a previsão opera em regime inicial.
\end{itemize} \\ \hline

Requisitos Não Funcionais Associados & 
\begin{itemize}[nosep, leftmargin=0.5cm]
    \item \textbf{Desempenho (Tempo de Resposta):} O processamento completo (recepção, inferência neural das três tarefas e gravação) deve ocorrer em tempo inferior a 500 milissegundos por requisição.
    \item \textbf{Concorrência e Disponibilidade:} O carregamento do modelo neural em memória deve ser único e protegido contra condições de corrida (\emph{Thread-Safe}), suportando múltiplas estações transmitindo simultaneamente.
\end{itemize} \\ \hline
\end{longtable}

% ============================================================================
% CASO DE USO 04
% ============================================================================
\section*{2. Especificação do Caso de Uso UC04}

\begin{longtable}{|>{\bfseries}p{3.8cm}|p{12.4cm}|}
\hline
\multicolumn{2}{|c|}{\textbf{UC04 --- Consultar Séries Históricas e Previsões}} \\ \hline

Identificador & \texttt{UC04} \\ \hline

Nome & Consultar Séries Históricas e Previsões \\ \hline

Breve Descrição & Este caso de uso permite ao usuário visualizar em painel interativo as medições reais coletadas pelas estações do campus, compará-las com as previsões geradas para 15 minutos à frente e acompanhar a classificação atual e futura do Índice de Qualidade do Ar (IQAr). \\ \hline

Atores & 
\begin{minipage}[t]{\linewidth}
\vspace{0.05cm}
\begin{tabular}{@{}p{4.2cm}p{7.8cm}@{}}
\textbf{Ator} & \textbf{Como participa do Caso de Uso} \\ \hline
\emph{Gestor Ambiental / Usuário} (Ator Primário) & Solicita a visualização dos dados ambientais, seleciona a estação sensora ou período desejado e analisa os gráficos e indicadores de qualidade do ar. \\
\end{tabular}
\vspace{0.1cm}
\end{minipage} \\ \hline

Prioridade & \textbf{Essencial (Alta)} \\ \hline

Pré-condições & 
\begin{enumerate}[nosep, leftmargin=0.5cm]
    \item Deve existir ao menos uma estação sensora cadastrada no sistema.
\end{enumerate} \\ \hline

Pós-condições & 
\begin{enumerate}[nosep, leftmargin=0.5cm]
    \item O painel analítico apresenta ao usuário as curvas temporais dos sete gases, as projeções futuras de 15 minutos, o indicador normativo CONAMA/EPA e o estado operacional dos sensores (nenhum dado é alterado no sistema).
\end{enumerate} \\ \hline

Fluxo de Eventos Principal (Caminho Feliz) & 
\begin{enumerate}[nosep, leftmargin=0.5cm]
    \item O \textbf{Gestor Ambiental / Usuário} acessa o painel de monitoramento da qualidade do ar no sistema.
    \item O \textbf{Sistema} recupera a lista de estações sensoras do campus e exibe o resumo geral do último Índice de Qualidade do Ar (IQAr) e o status de integridade de cada estação.
    \item O \textbf{Gestor Ambiental / Usuário} seleciona uma estação sensora específica (por exemplo, \emph{Laboratório de Engenharia}) e o intervalo de tempo ou composto gasoso de interesse.
    \item O \textbf{Sistema} consulta as séries históricas de medições reais e as respectivas previsões de 15 minutos geradas pelo Agente Inteligente para o filtro selecionado.
    \item O \textbf{Sistema} apresenta os gráficos comparativos (valor medido versus valor previsto), as probabilidades de cada classe normativa (\emph{Boa}, \emph{Moderada}, \emph{Ruim}) e o erro de reconstrução do sensor, encerrando o caso de uso.
\end{enumerate} \\ \hline

Fluxos Secundários (Alternativos e de Exceção) & 
\begin{itemize}[nosep, leftmargin=0.5cm]
    \item \textbf{[Exceção no Passo 4 --- Sem Leituras no Período Selecionado]:} Caso a estação escolhida não possua registros no período filtrado (por exemplo, equipamento desligado para manutenção), o \textbf{Sistema} exibe uma mensagem informativa indicando ausência de telemetria no intervalo e mantém visível a data/hora da última comunicação registrada.
\end{itemize} \\ \hline

Requisitos Não Funcionais Associados & 
\begin{itemize}[nosep, leftmargin=0.5cm]
    \item \textbf{Desempenho:} As consultas de séries históricas devem responder em no máximo 2 segundos, utilizando indexação temporal sobre a data e hora de coleta.
    \item \textbf{Usabilidade:} As faixas de qualidade do ar (\emph{Boa}, \emph{Moderada} e \emph{Ruim}) devem ser destacadas visualmente segundo a codificação de cores oficial da norma CONAMA nº 491/2018.
\end{itemize} \\ \hline
\end{longtable}

% ============================================================================
% CASO DE USO 05
% ============================================================================
\section*{3. Especificação do Caso de Uso UC05}

\begin{longtable}{|>{\bfseries}p{3.8cm}|p{12.4cm}|}
\hline
\multicolumn{2}{|c|}{\textbf{UC05 --- Gerenciar Alertas Ativos}} \\ \hline

Identificador & \texttt{UC05} \\ \hline

Nome & Gerenciar Alertas Ativos \\ \hline

Breve Descrição & Este caso de uso permite ao gestor ambiental consultar os alertas ativos gerados pelo Agente Inteligente (tanto por previsão de poluição crítica quanto por falha física nos sensores), inspecionar o diagnóstico detalhado e registrar a resolução do incidente após a tomada de providências no local. \\ \hline

Atores & 
\begin{minipage}[t]{\linewidth}
\vspace{0.05cm}
\begin{tabular}{@{}p{4.2cm}p{7.8cm}@{}}
\textbf{Ator} & \textbf{Como participa do Caso de Uso} \\ \hline
\emph{Gestor Ambiental / Usuário} (Ator Primário) & Consulta a fila de alertas pendentes, verifica a causa do alerta (ambiental ou falha de hardware) e marca o alerta como resolvido após intervir no ambiente ou acionar manutenção. \\
\end{tabular}
\vspace{0.1cm}
\end{minipage} \\ \hline

Prioridade & \textbf{Importante (Média/Alta)} \\ \hline

Pré-condições & 
\begin{enumerate}[nosep, leftmargin=0.5cm]
    \item O usuário deve estar com acesso ativo ao painel de gestão de alertas.
    \item Deve existir ao menos um alerta registrado pelo sistema.
\end{enumerate} \\ \hline

Pós-condições & 
\begin{enumerate}[nosep, leftmargin=0.5cm]
    \item O status do alerta selecionado é atualizado para \emph{``Resolvido''}, registrando-se a data/hora do atendimento e liberando o sistema para novas notificações daquela estação caso o problema volte a ocorrer.
\end{enumerate} \\ \hline

Fluxo de Eventos Principal (Caminho Feliz) & 
\begin{enumerate}[nosep, leftmargin=0.5cm]
    \item O \textbf{Gestor Ambiental / Usuário} solicita a listagem de alertas ativos no sistema.
    \item O \textbf{Sistema} apresenta a lista de alertas não resolvidos ordenados por severidade e data/hora de ocorrência, exibindo a estação sensora afetada, o tipo de ocorrência (\emph{Poluição Prevista} ou \emph{Falha no Sensor}) e a mensagem de diagnóstico.
    \item O \textbf{Gestor Ambiental / Usuário} seleciona um alerta específico para visualizar os detalhes da leitura e da previsão que motivaram o disparo.
    \item O \textbf{Sistema} exibe os valores registrados no instante do alerta e o histórico recente da estação.
    \item Após adotar a medida cabível (ex.: acionar ventilação da sala ou encaminhar o sensor para manutenção), o \textbf{Gestor Ambiental / Usuário} aciona a opção de confirmar a resolução do alerta.
    \item O \textbf{Sistema} atualiza a situação do alerta para resolvido, registra o instante de encerramento e atualiza a lista de pendências na tela.
\end{enumerate} \\ \hline

Fluxos Secundários (Alternativos e de Exceção) & 
\begin{itemize}[nosep, leftmargin=0.5cm]
    \item \textbf{[Fluxo Alternativo no Passo 2 --- Consulta de Histórico de Alertas Resolvidos]:} O \textbf{Gestor Ambiental / Usuário} pode optar por filtrar alertas já encerrados anteriormente para fins de auditoria e relatórios de segurança do campus. O \textbf{Sistema} exibe o histórico completo sem habilitar a ação de resolução.
\end{itemize} \\ \hline

Requisitos Não Funcionais Associados & 
\begin{itemize}[nosep, leftmargin=0.5cm]
    \item \textbf{Confiabilidade (Controle de Tempestade de Alertas):} Enquanto um alerta permanecer ativo dentro da janela de \emph{cooldown} de 15 minutos, o sistema não deve duplicar notificações idênticas para a mesma estação.
\end{itemize} \\ \hline
\end{longtable}

% ============================================================================
% CASO DE USO 06
% ============================================================================
\section*{4. Especificação do Caso de Uso UC06}

\begin{longtable}{|>{\bfseries}p{3.8cm}|p{12.4cm}|}
\hline
\multicolumn{2}{|c|}{\textbf{UC06 --- Atualizar Modelo Neural (Hot - Reloading)}} \\ \hline

Identificador & \texttt{UC06} \\ \hline

Nome & Atualizar Modelo Neural (Hot - Reloading) \\ \hline

Breve Descrição & Este caso de uso permite ao administrador substituir e recarregar em tempo de execução os arquivos do modelo de Inteligência Artificial e seus normalizadores após um novo ciclo de treinamento, sem interromper o funcionamento do servidor nem perder pacotes de telemetria das estações em campo. \\ \hline

Atores & 
\begin{minipage}[t]{\linewidth}
\vspace{0.05cm}
\begin{tabular}{@{}p{4.2cm}p{7.8cm}@{}}
\textbf{Ator} & \textbf{Como participa do Caso de Uso} \\ \hline
\emph{Administrador do Sistema} (Ator Primário) & Disponibiliza os novos artefatos treinados da rede neural e comanda a recarga a quente (\emph{Hot-Reloading}) no sistema. \\
\end{tabular}
\vspace{0.1cm}
\end{minipage} \\ \hline

Prioridade & \textbf{Importante (Média)} \\ \hline

Pré-condições & 
\begin{enumerate}[nosep, leftmargin=0.5cm]
    \item O \textbf{Administrador do Sistema} deve estar autenticado com perfil administrativo.
    \item Os novos arquivos do modelo neural (\texttt{.keras}) e transformadores de normalização (\texttt{.pkl}) devem estar disponíveis no diretório de artefatos do servidor.
\end{enumerate} \\ \hline

Pós-condições & 
\begin{enumerate}[nosep, leftmargin=0.5cm]
    \item O modelo neural em memória RAM é substituído atomicamente pela nova versão e passa a atender imediatamente às próximas requisições de inferência (\texttt{UC02}).
\end{enumerate} \\ \hline

Fluxo de Eventos Principal (Caminho Feliz) & 
\begin{enumerate}[nosep, leftmargin=0.5cm]
    \item O \textbf{Administrador do Sistema} solicita ao sistema a atualização em tempo de execução (\emph{Hot-Reloading}) do modelo do Agente Inteligente.
    \item O \textbf{Sistema} localiza os novos arquivos do modelo neural e dos normalizadores no repositório de artefatos e realiza o carregamento preliminar em espaço isolado de memória, sem interromper as inferências em andamento.
    \item O \textbf{Sistema} valida a estrutura das três saídas do novo modelo (reconstrução, previsão de 15 minutos e classificação do IQAr).
    \item O \textbf{Sistema} aplica bloqueio de exclusão mútua (\emph{lock}), substitui a referência do modelo antigo pelo novo modelo validado em memória e libera o bloqueio.
    \item O \textbf{Sistema} exibe mensagem de confirmação ao \textbf{Administrador do Sistema} informando a versão carregada e o horário da atualização.
\end{enumerate} \\ \hline

Fluxos Secundários (Alternativos e de Exceção) & 
\begin{itemize}[nosep, leftmargin=0.5cm]
    \item \textbf{[Exceção nos Passos 2 e 3 --- Arquivo Ausente, Corrompido ou Incompatível]:} Caso o novo arquivo do modelo esteja corrompido ou apresente formato incompatível, o \textbf{Sistema} aborta a substituição, preserva o modelo anterior funcionando normalmente em memória (\emph{graceful fallback}), registra o erro em log e notifica o \textbf{Administrador do Sistema} sobre a falha.
\end{itemize} \\ \hline

Requisitos Não Funcionais Associados & 
\begin{itemize}[nosep, leftmargin=0.5cm]
    \item \textbf{Disponibilidade (\emph{Zero-Downtime}):} A atualização do modelo não pode derrubar o serviço nem recusar pacotes enviados pelo ChirpStack durante a troca.
    \item \textbf{Segurança:} Operação restrita exclusivamente a usuários autenticados com permissão de Administrador do Sistema.
\end{itemize} \\ \hline
\end{longtable}

% ============================================================================
% CASO DE USO 07 (Com Subfluxos conforme exemplo da Vídeo-Aula)
% ============================================================================
\section*{5. Especificação do Caso de Uso UC07}

\begin{longtable}{|>{\bfseries}p{3.8cm}|p{12.4cm}|}
\hline
\multicolumn{2}{|c|}{\textbf{UC07 --- Gerenciar Nós Sensores (Sensor Node)}} \\ \hline

Identificador & \texttt{UC07} \\ \hline

Nome & Gerenciar Nós Sensores (Sensor Node) \\ \hline

Breve Descrição & Este caso de uso permite ao administrador manter o inventário das estações sensoras instaladas no campus universitário, abrangendo o cadastramento de novos dispositivos LoRaWAN, a atualização de localização/parâmetros e a desativação de nós sensores. \\ \hline

Atores & 
\begin{minipage}[t]{\linewidth}
\vspace{0.05cm}
\begin{tabular}{@{}p{4.2cm}p{7.8cm}@{}}
\textbf{Ator} & \textbf{Como participa do Caso de Uso} \\ \hline
\emph{Administrador do Sistema} (Ator Primário) & Informa os dados cadastrais da estação sensora (\texttt{DevEUI}, nome, localização física no campus e status operacional). \\
\end{tabular}
\vspace{0.1cm}
\end{minipage} \\ \hline

Prioridade & \textbf{Importante (Média)} \\ \hline

Pré-condições & 
\begin{enumerate}[nosep, leftmargin=0.5cm]
    \item O \textbf{Administrador do Sistema} deve estar autenticado no painel administrativo do sistema.
\end{enumerate} \\ \hline

Pós-condições & 
\begin{enumerate}[nosep, leftmargin=0.5cm]
    \item Os dados cadastrais da estação sensora são incluídos, atualizados ou desativados no banco de dados do sistema.
\end{enumerate} \\ \hline

Fluxo de Eventos Principal e Subfluxos & 
\begin{enumerate}[nosep, leftmargin=0.5cm]
    \item O \textbf{Administrador do Sistema} acessa a funcionalidade de gerenciamento de nós sensores.
    \item O \textbf{Sistema} exibe a relação das estações sensoras cadastradas no campus e disponibiliza as opções de \emph{Incluir Nó Sensor}, \emph{Alterar Informações do Nó Sensor} ou \emph{Desativar Nó Sensor}.
    \item O \textbf{Administrador do Sistema} seleciona a operação desejada.
    \item O \textbf{Sistema} executa o respectivo subfluxo selecionado:
    \begin{itemize}[nosep, leftmargin=0.5cm]
        \item \textbf{Subfluxo 4.1 --- Incluir Nó Sensor:}
        \begin{enumerate}[nosep, leftmargin=0.5cm]
            \item O \textbf{Administrador} informa o código identificador único do rádio (\texttt{DevEUI} de 16 caracteres hexadecimais), o nome da estação e o local de instalação (ex.: \emph{Bloco B --- Sala 204}).
            \item O \textbf{Sistema} valida que o \texttt{DevEUI} é único, registra a nova estação como ativa e confirma o cadastro.
        \end{enumerate}
        \item \textbf{Subfluxo 4.2 --- Alterar Informações do Nó Sensor:}
        \begin{enumerate}[nosep, leftmargin=0.5cm]
            \item O \textbf{Administrador} seleciona um nó sensor existente e edita sua descrição ou localização física após remanejamento no campus.
            \item O \textbf{Sistema} valida os dados, atualiza o registro mantendo o histórico de leituras preservado e confirma a alteração.
        \end{enumerate}
        \item \textbf{Subfluxo 4.3 --- Desativar Nó Sensor:}
        \begin{enumerate}[nosep, leftmargin=0.5cm]
            \item O \textbf{Administrador} seleciona uma estação recolhida para manutenção e solicita sua desativação.
            \item O \textbf{Sistema} marca o nó sensor como inativo (preservando todo o histórico científico já coletado) e encerra o subfluxo.
        \end{enumerate}
    \end{itemize}
\end{enumerate} \\ \hline

Fluxos Secundários (Alternativos e de Exceção) & 
\begin{itemize}[nosep, leftmargin=0.5cm]
    \item \textbf{[Exceção no Subfluxo 4.1 --- Código \texttt{DevEUI} Duplicado ou Inválido]:} Caso o \textbf{Administrador} informe um código \texttt{DevEUI} já cadastrado em outra estação ou fora do padrão hexadecimal de 16 dígitos, o \textbf{Sistema} rejeita o cadastro, exibe mensagem indicando o conflito e solicita a correção do valor.
\end{itemize} \\ \hline

Requisitos Não Funcionais Associados & 
\begin{itemize}[nosep, leftmargin=0.5cm]
    \item \textbf{Integridade Histórica:} A desativação de um nó sensor jamais deve apagar as séries históricas de poluentes e previsões já registradas para auditoria científica.
\end{itemize} \\ \hline
\end{longtable}

\end{document}
